\ysrc{65}
再举一个稍复杂些的例子：取 $M$ 为 $\mathcal O_E$ 上的分裂连通约化群，记 $M/\!/M$ 为其伴随商。则
\[
\mathcal R^c_{\Gamma,M}\cong M\times_{M/\!/M}\Big(\bigsqcup_x (M/\!/M)^{\wedge}_x\Big),
\]
其中 $x$ 取遍 $M/\!/M$ 的所有闭点。

\begin{rem}{2.51}
例 2.50 表明：定义 2.38 虽然对任意局部投射有限群 $\Gamma$ 都有意义，但在某些用途下它给出的未必是「最正确」的对象。也就是说，尽管 $\mathcal R^c_{\Gamma,M}$ 已经把 $\Gamma$ 的各种形变空间粘在一起，它一般仍不连通，并且带有形式方向。这个例子还提示，在有些情形下 $\mathcal R^c_{\Gamma,M}$ 的不同分支可以进一步粘合。例如 $\mathcal R^c_{\widehat{\mathbb Z},M}$ 的所有分支应当自然地粘成 $\mathcal R_{\mathbb Z,M}=M$。这是下文讨论的一般现象的一个特例（特别见命题 2.52）。再举一例：设 $F$ 是剩余特征为 $p$ 的非阿基米德局部域，$\Gamma_F$ 为其伽罗瓦群，$W_F$ 为其 Weil 群。若 $p\neq\ell$，考虑 $\mathcal R^c_{W_F,M}$ 要比考虑 $\mathcal R^c_{\Gamma_F,M}$ 更正确，下一节就会看到这一点。若 $p=\ell$，则连 $\mathcal R^c_{W_F,M}$ 也不够，这是 Emerton 告诉我们的；此时需要 [23] 中的构造。最后，当 $\Gamma$ 是 $\mathbb F_p$ 上一条光滑（仿射）代数曲线的 étale 基本群（且 $p\neq\ell$）时，我们也预期它的表示空间有更精细的构造。
\end{rem}

当 $\widehat\Gamma$ 是抽象群 $\Gamma$ 的投射有限完备化时（如引理 2.45 中那样），在若干温和的假设下，$\mathcal R_{\Gamma,M}$ 会把 $\mathcal R^c_{\widehat\Gamma,M}$ 的不同分支粘合起来（如分解 (2.36) 所示）。

\begin{prop}{2.52}
设 $\Gamma$ 是 $FP_\infty(k)$ 型的有限生成群，使得映射 $\Gamma\to\widehat\Gamma$ 对每个有限 $\mathbb F_\ell\Gamma$-模 $V$ 都诱导群上同调的同构 $H^i_{cts}(\widehat\Gamma,V)\cong H^i(\Gamma,V)$。则自然态射 $\mathcal R^c_{\widehat\Gamma,M}\to\mathcal R_{\Gamma,M}$ 诱导同构
\[
\mathcal R^c_{\widehat\Gamma,M}\cong \mathcal R_{\Gamma,M}\times_{\mathcal R_{\Gamma,M^{\bullet}/\!/H}}\Big(\bigsqcup_x\big(\mathcal R_{\Gamma,M^{\bullet}/\!/H}\big)^{\wedge}_x\Big),
\]
其中 $x$ 取遍 $\mathcal R_{\Gamma,M^{\bullet}/\!/H}$ 在 $\kappa_E$ 上的所有闭点，而 $\big(\mathcal R_{\Gamma,M^{\bullet}/\!/H}\big)^{\wedge}_x$ 是 $\mathcal R_{\Gamma,M^{\bullet}/\!/H}$ 在 $x$ 处的形式完备化。
\end{prop}

\begin{proof}
我们只在经典模问题的层面给出证明。与命题 2.42 类似的论证表明，它在导出层面也是同构。

由命题 2.31，$\mathcal R^c_{\widehat\Gamma,M}\to\mathcal R_{\Gamma,M}$ 显然穿过态射
$\mathcal R^c_{\widehat\Gamma,M}\to \mathcal R_{\Gamma,M}\times_{\mathcal R_{\Gamma,M^{\bullet}/\!/H}}\big(\bigsqcup_x(\mathcal R_{\Gamma,M^{\bullet}/\!/H})^{\wedge}_x\big)$。只需构造逆映射。

设 $\rho:\Gamma\to M(A)$ 是同态，其中 $A$ 是 $\mathcal O_E$ 上有限型的经典代数，使得复合态射
$\Spec A\to\mathcal R_{\Gamma,M}\to\mathcal R_{\Gamma,M^{\bullet}/\!/H}$ 把拓扑空间 $|\Spec A|$ 映入 $x$。只需证明 $\rho$ 穿过 $\Gamma$ 的一个有限商。可取忠实嵌入 $M\hookrightarrow \mathrm{GL}_m$ 并设 $M=\mathrm{GL}_m$。按假设，映射 $\Lambda[\mathcal R_{\Gamma,\mathrm{GL}_m^{\bullet}}/\!/\mathrm{GL}_m]^{\mathrm{GL}_m}\to A$ 的像 $B$ 是 Artin 局部环。

注意对每个 $\gamma\in\Gamma$，$\rho(\gamma):A^m\to A^m$ 的特征多项式 $\mathrm{Char}(\rho(\gamma),t)=\det(t-\rho(\gamma))$ 属于 $B[t]$。下面的论证是 [22, 2.8-2.10] 的轻微变体。

先设 $A$ 约化，从而是有限型 $\kappa_E$-代数。则 $B$ 是 $\kappa_E$ 的有限扩张。我们知道存在 $B$ 的有限扩张 $\kappa$ 以及完全可约表示 $\rho':\Gamma\to\mathrm{GL}_m(\kappa)$，使得对每个 $\gamma\in\Gamma$ 都有 $\mathrm{Char}(\rho'(\gamma),t)=\mathrm{Char}(\rho(\gamma),t)$。特别地，存在有限指数子群 $\Gamma_1\subset\Gamma$，使得 $\mathrm{Char}(\rho(\gamma),t)=(t-1)^m$。

\ysrc{66}
把 $A$ 换成其商环并作共轭，可设对每个 $\gamma\in\Gamma_1$，$\rho(\gamma)$ 都是严格上三角的。注意系数在 $A$ 中的严格上三角矩阵群是幂零群，其指数为 $\ell$ 的某个幂。按假设 $H^1(\Gamma_1,\mathbb F_\ell)$ 是有限维 $\mathbb F_\ell$-向量空间。于是存在有限指数子群 $\Gamma_2\subset\Gamma_1$，使得 $\rho|_{\Gamma_2}$ 平凡。

对一般的有限型 $\mathcal O_E$-代数 $A$（其中 $\ell$ 幂零），设 $A_{red}$ 为它模去幂零根所得的商。令 $\Gamma_2$ 为 $\Gamma\to\mathrm{GL}_m(A)\to\mathrm{GL}_m(A_{red})$ 的核，它在 $\Gamma$ 中有限指数。由于核 $\mathrm{GL}_m(A)\to\mathrm{GL}_m(A_{red})$ 是幂零群，指数为 $\ell$ 的某个幂，而 $H^1(\Gamma_2,\mathbb F_\ell)$ 有限维，故存在有限指数子群 $\Gamma_3\subset\Gamma_2$，使得 $\rho|_{\Gamma_3}$ 平凡。
\end{proof}

本小节的最后一个话题是把模空间 $\mathcal R^c_{\Gamma,M}$ 从 $\Spf\mathcal O_E$ 推广到 $\Spec\mathcal O_E$。当然，若 $\Gamma$ 恰是某个抽象群 $\Gamma_0$ 的投射有限完备化（如命题 2.52 中那样），这样的推广可由 $\mathcal R_{\Gamma_0,M}$ 给出。这正是我们在构造局部 Langlands 参数模空间时（在 $\ell\neq p$ 情形）要采用的办法。然而并非每个 $\Gamma$ 都以这种方式出现，而且即使如此，一般也没有 $\Gamma_0$ 的典范选取。因此最好有一个更直接的构造。由于一般地 $\mathcal R^c_{\Gamma,M}$ 带有非平凡的形式方向，这样的推广在一般情形下大概应当是解析性质的。不过，对下一节所考虑的具体情形，下述办法已经足够。

思路是把引理 2.36 中对 $\mathcal O_{E,r}$-模/代数定义的 $C_{cts}(S,-)$ 推广为对满足类似性质的 $\mathcal O_E$-模/代数都有定义的函子。此后几乎所有其余构造都可照搬。

设 $\mathbf{Mod}^{\heartsuit,f.g.}_{\mathcal O_{E,r}}$ 表示有限 $\mathcal O_{E,r}$-模的阿贝尔范畴。从它到有限集范畴 $\mathbf{Sets}_f$ 的自然遗忘函子忠实、保守、保持有限积，并且是松对称幺半的（这里 $\mathbf{Sets}_f$ 赋以 Descartes 对称幺半结构）。它诱导自然函子\footnote{考虑这种函子的想法来自 Peter Scholze，他用凝聚数学（condensed mathematics）发展了一套连续表示模空间的方法。我们这里的办法没有用到凝聚数学，但很可能与 Scholze 的本质相同。}
\begin{equation*}\tag{2.37}
\mathbf{Mod}^{\heartsuit}_{\mathcal O_E}=\Ind\varprojlim_r\mathbf{Mod}^{\heartsuit,f.g.}_{\mathcal O_{E,r}}\longrightarrow \Ind\mathrm{Pro}(\mathbf{Sets}_f),
\end{equation*}
并满足类似性质，其中 $\Ind\mathrm{Pro}(\mathbf{Sets}_f)$ 表示投射有限集范畴的 ind-完备化。注意投射有限集的不交并 $S$ 也可以看作 $\Ind\mathrm{Pro}(\mathbf{Sets}_f)$ 中的对象。

\begin{lem}{2.53}
设 $S$ 是投射有限集的不交并，看作 ind-投射有限集。则
$\mathbf{Mod}^{\heartsuit}_{\mathcal O_E}\to\mathbf{Mod}^{\heartsuit}_{\mathcal O_E}$，$C_{cts}(S,V)=\Map_{\Ind\mathrm{Pro}(\mathbf{Sets}_f)}(S,V)$
满足与引理 2.36 中相同的性质，从而延拓为 t-正合函子
\begin{equation*}\tag{2.38}
C_{cts}(S,-):\mathbf{Mod}_{\mathcal O_E}\longrightarrow\mathbf{Mod}_{\mathcal O_E},
\end{equation*}

\ysrc{67}
它又提升为幂完备（nilcomplete）函子
\begin{equation*}\tag{2.39}
C_{cts}(S,-):\mathbf{CAlg}_{\mathcal O_E}\longrightarrow\mathbf{CAlg}_{\mathcal O_E},
\end{equation*}
并保持有限极限。若 $S$ 投射有限，则 (2.38) 保持所有余极限，而 (2.39) 保持筛（sifted）余极限。
\end{lem}

\begin{proof}
关于模的第一部分，用引理 2.36 的论证，可归约为：当 $S$ 投射有限时，对有限 $\mathcal O_E$-模的满射 $M\to M''$ 证明 $C_{cts}(S,M)\to C_{cts}(S,M'')$ 满。由于每个有限 $\mathcal O_E$-模都是有限自由模与有限挠模的直和，这一点也就清楚了。由于 (2.37) 是松幺半的，当 $A$ 是 $\mathcal O_E$-代数时 $C_{cts}(S,A)$ 也是。其余部分的论证与引理 2.36 相同。
\end{proof}

\begin{rem}{2.54}
注意由 Stone 对偶，全忠实函子 $\mathrm{Pro}(\mathbf{Sets}_f)\to\mathbf{Top}$ 诱导全忠实函子 $\Ind\mathrm{Pro}(\mathbf{Sets}_f)\to\mathbf{Top}$。结合 (2.37)，这给每个 $\mathcal O_E$-模赋予一个拓扑，我们称之为 ind-$\varpi$-进拓扑。具体地说，对 $\mathcal O_E$-模 $V$，它是 $V$ 上使得每个有限生成子模 $U\subset V$ 的子空间拓扑都与 $\varpi$-进拓扑一致的最细拓扑。一般地，$V$ 上的 ind-$\varpi$-进拓扑强于 $V$ 上的某些其他方便的拓扑。例如若 $V$ 是 $\varpi$-进分离的 $\mathcal O_E$-模，则 $V$ 上的 ind-$\varpi$-进拓扑通常严格细于 $\varpi$-进拓扑。同样，对代数扩张 $F/E$，$F$（看作 $\mathcal O_E$-模）上的 ind-$\varpi$-进拓扑严格细于 $F$ 上通常的 $\varpi$-进拓扑，除非 $[F:E]<\infty$。注意若 $V$ 对某个 $r$ 是 $\mathcal O_{E,r}$-模，则 $V$ 上的 ind-$\varpi$-进拓扑是离散的。

还有一点需要提醒。由于函子 $\Ind\mathrm{Pro}(\mathbf{Sets}_f)\to\mathbf{Top}$ 一般不保持有限积，复合函子 $\mathbf{Mod}^{\heartsuit}_{\mathcal O_E}\to\mathbf{Top}$ 不是松对称幺半的，因此赋以 ind-$\varpi$-进拓扑的经典 $\mathcal O_E$-代数 $A$ 未必是通常意义下的拓扑代数。弥补这一点的办法是注意到 $\Ind\mathrm{Pro}(\mathbf{Sets}_f)\to\mathbf{Top}$ 实际上穿过 $\Ind\mathrm{Pro}(\mathbf{Sets}_f)\to\mathbf{CG}$，其中 $\mathbf{CG}\subset\mathbf{Top}$ 是紧生成空间的全子范畴，而由此得到的函子保持有限积。
\end{rem}

\begin{rem}{2.55}
如同 $\mathcal O_{E,r}$ 上的情形，把 $S$ 看作离散集，有 $C_{cts}(S,-)\to C(S,-)$。若 $A$ 是经典的，则 $C_{cts}(S,A)\to C(S,A)$ 是单射。此外注意若 $V\in\mathbf{Mod}^{\heartsuit,f.g.}_{\mathcal O_E}$，则 $C_{cts}(S,V)=\varprojlim_r C_{cts}(S,V/\varpi^r)$。
\end{rem}

现在给定引理 2.53 中的 $C_{cts}(S,-)$，我们可以如下推广定义 2.38。

\begin{defn}{2.56}
把 $\mathcal O_E$ 上取值于 $M$ 的强连续表示空间定义为
\[
\mathcal R^{sc}_{\Gamma,M}:\mathbf{CAlg}_{\mathcal O_E}\to\mathbf{Spc},\qquad
A\mapsto \Map_{\mathbf{CAlg}^{\Delta}_{\mathcal O_E}}\big(\mathcal O_E[M^{\bullet}],C_{cts}(\Gamma^{\bullet},A)\big).
\]
类似地，把表示叠 $\mathcal R^{sc}_{\Gamma,M/H}$ 定义为 $\mathcal R^{sc}_{\Gamma,M}$ 对 $H$ 的商。存在与 (2.19) 类似的单纯图表（把 $\mathcal R_{\Gamma,M}$ 换成 $\mathcal R^{sc}_{\Gamma,M}$），于是我们把表示叠 $\mathcal R^{sc}_{\Gamma,M/H}$ 定义为它在 $\mathbf{Shv}(\mathbf{CAlg}_{\mathcal O_E})$ 中的几何实现。类似地，把 $\Spec\mathcal O_E$ 上的伪表示空间定义为
\[
\mathcal R^{sc}_{\Gamma,M^{\bullet}/\!/H}:\mathbf{CAlg}_{\mathcal O_E}\to\mathbf{Spc},\qquad
A\mapsto \Map_{\mathbf{CAlg}^{\mathrm{FFM}}_{\mathcal O_E}}\big(\mathcal O_E[M^{\bullet}]^{H},C_{cts}(\Gamma^{\bullet},A)\big).
\]
\end{defn}

\ysrc{68}
按定义，$\mathcal R^{sc}_{\Gamma,M}$ 在 $\Spf\mathcal O_E$ 上的限制就是 $\mathcal R^c_{\Gamma,M}$。如前，有自然态射
\[
\mathcal R^{sc}_{\Gamma,M}\to\mathcal R_{\Gamma,M},\qquad
\mathcal R^{sc}_{\Gamma,M/H}\to\mathcal R_{\Gamma,M/H}
\]
（都定义在 $\mathcal O_E$ 上），其中 $\mathcal R_{\Gamma,M}$ 与 $\mathcal R_{\Gamma,M/H}$ 中的 $\Gamma$ 看作抽象群。若 $A$ 是经典的，则由 $\mathcal R^{sc}_{\Gamma,M}(A)\to\mathcal R_{\Gamma,M}(A)$ 诱导的映射是单射，其像由那些 $\rho:\Gamma\to M(A)$ 组成：对每个 $f\in\mathcal O_E[M]$，$f\circ\rho:\Gamma\to A$ 都连续，其中 $A$ 赋以 ind-$\varpi$-进拓扑。由于 $A$ 上的 ind-$\varpi$-进拓扑一般强于其他方便的拓扑（见注 2.54），我们称这样的 $\rho$ 为强连续表示（strongly continuous representation）。这正是 $\mathcal R^{sc}_{\Gamma,M}$ 这一术语的由来。

下面这个简单的观察在后续许多讨论中都很重要。

\begin{lem}{2.57}
设 $\Gamma$ 投射有限，$A$ 是经典 $\mathcal O_E$-代数。则 $\rho:\Gamma\to\mathrm{GL}_m(A)$ 属于 $\mathcal R^{sc}_{\Gamma,\mathrm{GL}_m}(A)$，当且仅当 $A^m=\bigcup_i V_i$ 是有限 $\mathcal O_E$-模 $V_i$ 的并，每个 $V_i$ 都是 $\Gamma$-稳定的，且 $\Gamma$ 在 $V_i$ 上的作用连续。
\end{lem}

\begin{proof}
事实上，若记 $\rho(\gamma)$ 的 $(i,j)$-元为 $a_{ij}(\gamma)$，则 $\Gamma\to A$，$\gamma\mapsto a_{ij}(\gamma)$ 是 $\Ind\mathrm{Pro}(\mathbf{Sets}_f)$ 中的映射，故其像含于 $A$ 的一个有限生成 $\mathcal O_E$-子模。因此对每个 $v\in A^m$，$\rho(\Gamma)v$ 含于 $A^m$ 的一个在 $\mathcal O_E$ 上有限的 $\Gamma$-子模 $V$，且 $\Gamma$ 在 $V$ 上的作用连续。反之，若 $A^m$ 是引理中那样的 $\Gamma$-子模 $V_i$ 之并，则 $a_{ij}:\Gamma\to A$ 取值于 $A$ 的有限生成 $\mathcal O_E$-子模，且由此得到的映射连续。于是 $\rho$ 强连续。
\end{proof}

\begin{rem}{2.58}
利用上述引理可以证明 ${}^{cl}\mathcal R^{sc}_{\Gamma,M}$ 可由 ind-仿射概形表示。由于我们不用这一事实，略去证明。
\end{rem}

现在对 $\rho\in\mathcal R^{sc}(\Gamma,M)(A)$ 以及 $M$ 在有限自由 $\mathcal O_E$-模 $W$ 上的代数表示 $W$，同样有 $W_\rho=W\otimes A$，带有 $\Gamma$ 的强连续作用（由 $C_{cts}(\Gamma^{\bullet},A)$ 上方的余单纯模 $C_{cts}(\Gamma^{\bullet},W_\rho)$ 编码，如注 2.13(1)）。设 $C^*_{cts}(\Gamma,W_\rho)$ 为 $C_{cts}(\Gamma^{\bullet},W_\rho)$ 的全化（在 $\mathbf{Mod}_A$ 中）。鉴于注 2.16，我们把这个上链复形称为 $\Gamma$ 以 $W_\rho$ 为系数的连续群上同调。类似地有约化版本 $\overline{C}^*_{cts}(\Gamma,W_\rho)[1]$。

若 $A$ 经典且 $\Gamma$ 投射有限，则由引理 2.57 可写 $W_\rho=\bigcup_i V_i$，其中每个 $V_i$ 是 $\Gamma$ 在有限 $\mathcal O_E$-模上的连续表示。由于 $S$ 投射有限时 $C_{cts}(S,-)$ 与滤余极限交换，有
\begin{equation*}\tag{2.40}
C^*_{cts}(\Gamma,W_\rho)=\varinjlim_i C^*_{cts}(\Gamma,V_i),
\end{equation*}
其中 $C^*_{cts}(\Gamma,V_i)$ 是 $\Gamma$ 以连续 $\Gamma$-模 $V_i$ 为系数的通常连续群上同调。

下面的命题总结了 $\mathcal R^{sc}_{\Gamma,M}$ 的无穷小几何，它是关于 $\mathcal R^c_{\Gamma,M}$ 的相应陈述的方向推广\tn{底本此处作「direction generalization」，疑为「direct generalization」（直接推广）之误。}。

\begin{prop}{2.59}
函子 $\mathcal R^{sc}_{\Gamma,M}:\mathbf{CAlg}_{\mathcal O_E}\to\mathbf{Spc}$ 幂完备且保持有限积。设 $\rho\in\mathcal R^{sc}_{\Gamma,M}(A)$，$A$ 截断。则 $T_\rho\mathcal R^{sc}_{\Gamma,M}(A)\cong \overline{C}^*_{cts}(\Gamma,\mathrm{Ad}_\rho)[1]$。若 $\Gamma$ 投射有限，则对每个 $m$，$\mathcal R^{sc}_{\Gamma,M}$ 到函子 ${}^{\leq m}\mathbf{CAlg}_{\mathcal O_E}\to\mathbf{Spc}$ 的限制与滤余极限交换。
\end{prop}

\ysrc{69}
本小节以一个关于投射有限群的强连续表示的剩余伪表示之常性的结果结束。于是设 $\Gamma$ 投射有限，$(M,H)$ 如命题 2.31。首先，如 [7, 4.8] 所述，对每个连续表示 $\rho:\Gamma\to M(E')$（$E'/E$ 为有限扩张），$\mathrm{Tr}\,\rho$ 的伪表示取 $\mathcal O_{E'}$-值，故它模 $\varpi'$ 的约化给出一个良定义的 $\kappa_{E'}$-值伪表示，记作 $\overline{\mathrm{Tr}}\,\rho$。为统一记号，若 $\rho:\Gamma\to M(\kappa')$ 连续且 $\kappa'/\kappa_E$ 有限，也把 $\mathrm{Tr}\,\rho$ 记作 $\overline{\mathrm{Tr}}\,\rho$。

\begin{lem}{2.60}
设 $A$ 是有限生成 $\mathcal O_E$-代数，$\Spec A$ 连通，$\rho:\Gamma\to M(A)$ 是强连续表示。对每个剩余域在 $\kappa_E$ 上有限或在 $E$ 上有限的点 $x\in\Spec A$，记 $\rho_x$ 为相应的连续表示。则 $x\mapsto\overline{\mathrm{Tr}}\,\rho_x$ 是常值。
\end{lem}

\begin{proof}
若对某个 $n$ 有 $\varpi^nA=0$，这由命题 2.31 得到。今设 $A[\varpi^{-1}]$ 非空。设 $\Spec B\subset\Spec A[\varpi^{-1}]$ 是连通分支。设 $B_0$ 是 $B$ 中由所有 $f(\rho(\gamma_1,\ldots,\gamma_n))$（$n\geq 1$，$f\in E[M^n]^H$，$(\gamma_i)\in\Gamma^n$）生成的子环。由于 FFM-代数 $E[M^{\bullet}]^H$ 有限生成（[76, thm. 9]）且 $\rho$ 强连续，$B_0$ 在 $E$ 上有限生成。由于 $\Spec B_0$ 的每个闭点实际上都定义在 $\mathcal O_E$ 的某个有限扩张上，$B_0$ 本身必在 $E$ 上有限。由于 $\Spec B_0$ 连通，它只有一个点。故 $\mathrm{Tr}(\rho_x)$ 为常值。最后，显然若 $\rho_x:\Gamma\to M(E')$ 来自 $\rho_x:\Gamma\to M(\mathcal O_{E'})$，则 $\mathrm{Tr}\,\rho=\mathrm{Tr}\,\bar\rho=\overline{\mathrm{Tr}}\,\rho$。引理是上述事实的组合。
\end{proof}

\section{算术 Langlands 参数的叠}

本节把上一节的构造用于理解 Langlands 参数的模空间。局部情形（$\ell\neq p$）相对而言已经比较清楚，将在 §3.1–3.3 讨论。整体域情形能说的少得多；不过在整体函数域情形，我们仍能在 §3.4 构造该模空间。

先回忆 Buzzard–Gee [12] 引入的 C-群，它遵循 [87, §1.1] 的构造。设 $G$ 是域 $F$ 上的连通约化群。记 $\Gamma_F$ 为 $F$ 的伽罗瓦群，$\hatG$ 为 $G$ 的对偶群，看作 $\Z$ 上的群概形。它带有固定数据（pinning）$(\widehat{B},\widehat{T},\hat{e})$，以及经由同态 $\xi:\Gamma_F\to\Aut(\hatG,\widehat{B},\widehat{T},\hat{e})$ 的 $\Gamma_F$-作用。设 $\hatG_{ad}$ 为 $\hatG$ 的伴随群，$\rho_{ad}:\mathbb G_m\to\hatG_{ad}$ 为 $\hatG$ 的正余根之半和给出的余特征标。设 $\mathrm{pr}:\Gamma_F\to\Gamma_{\wt F/F}$ 为 $\Gamma_F$ 对 $\ker\xi$ 的有限商。令
\[
{}^cG:=\hatG\rtimes\big(\mathbb G_m\times\Gamma_{\wt F/F}\big)
\]
为 $G$ 的 C-群，看作 $\Z$ 上的群概形，其中 $\mathbb G_m$ 经由同态 $\mathbb G_m\xrightarrow{\rho_{ad}}\hatG_{ad}\subset\Aut(\hatG)$ 作用于 $\hatG$，而 $\Gamma_{\wt F/F}$ 经由 $\xi$ 作用。记 $d:{}^cG\to\mathbb G_m\times\Gamma_{\wt F/F}$ 为自然投影。

\begin{rem}{3.1}
若 $F$ 是剩余域为 $\mathbb F_q$ 的局部域，或是常数域为 $\mathbb F_q$ 的整体函数域，则在选定 $q^{1/2}$ 之后，${}^cG$ 与 ${}^LG\times\mathbb G_m$ 在 $\Z[q^{\pm1/2}]$ 上同构，其中 ${}^LG=\hatG\rtimes\Gamma_{\wt F/F}$ 是 $G$ 的通常 Langlands 对偶群。因此在下面大多数讨论中可以把 ${}^cG$ 换成 ${}^LG$（作小的修改）。然而在我们的表述中更愿意用 C-群而不用 L-群。一方面，它更典范（在考虑数域的局部-整体相容性时将是必要的）。另一方面，当 $\wt F\neq F$ 时，用 L-群似乎并不能把表述简化多少。
\end{rem}

\ysrc{70}
反过来，若余特征标 $\rho_{ad}$ 可以提升为 $\Gamma_{\wt F/F}$-不变的余特征标 $\tilde\rho:\mathbb G_m\to\hatG$，则下面的讨论中也可以用 ${}^LG$ 代替 ${}^cG$。例如 $G=\mathrm{GL}_n$ 或奇数维酉群时就是如此。见 [87, Example 2]。

\subsection{局部 Langlands 参数的叠}

下面两个小节讨论以 $p$ 可逆的环为基、对局部域 $F$（剩余特征为 $p$）上的连通约化群 $G$ 的局部 Langlands 参数叠。本小节的部分结果也由 Dat–Helm–Kurinczuk–Moss [19] 得到，并由 Fargues–Scholze [27] 独立得到，有时所用方法不同。

记 $\kappa_F$ 为剩余域，$\#\kappa_F=p^r$。设 $\Gamma_F$ 为 $F$ 的伽罗瓦群。设 $P_F\subset I_F\subset\Gamma_F$ 分别为野惯性群与惯性群，对应于伽罗瓦扩张 $F^t\supset F^{ur}\supset F$。回忆驯顺惯性群
\[
I^t_F:=I_F/P_F\cong\prod_{\ell\neq p}\Z_\ell(1)=:\hat\Z^p(1)
\]
是素于 $p$ 的，而 $P_F$ 是 pro-$p$ 群。于是 $\Gamma^t_F:=\Gamma_{F^t}\cong\Gamma_F/P_F$ 落在如下短正合列中
\[
1\to I^t_F\to\Gamma^t_F\to\hat\Z\to1.
\]
设 $W_F\subset\Gamma_F$ 为 $F$ 的 Weil 群。我们规范化映射
\begin{equation*}\tag{3.1}
\|\cdot\|:W_F\to\Z
\end{equation*}
使它在 $I_F$ 上平凡，且对算术 Frobenius 的某个提升 $\Phi$ 有 $\|\Phi\|=1$。类似地，有驯顺 Weil 群 $W^t_F:=W_F/P_F$，它是 $\Z$ 被 $I^t_F$ 的扩张。令
\[
\chi=(q^{-\|\cdot\|},\mathrm{pr}):W_F\to\Z[1/p]^\times\times\Gamma_{\wt F/F}.
\]
注意 $q^{-\|\cdot\|}$ 是 $\Gamma_F$ 的逆分圆特征标在 $W_F$ 上的限制。

局部 Langlands 参数的模有好几种版本。首先固定素数 $\ell\neq p$。有 $W_F$ 在 $\Spf\Z_\ell$ 上的连续表示模 $\mathcal R^c_{W_F,{}^cG}$（定义 2.38）。同态 $d:{}^cG\to\mathbb G_m\times\Gamma_{\wt F/F}$ 诱导态射 $\mathcal R^c_{W_F,{}^cG}\to\mathcal R^c_{W_F,\mathbb G_m\times\Gamma_{\wt F/F}}$。可以把 $\chi$ 看作 $\mathcal R^c_{W_F,\mathbb G_m\times\Gamma_{\wt F/F}}$ 的一个 $\Spf\Z_\ell$-点，并定义
\begin{equation*}\tag{3.2}
\Loc^{\wedge,\square}_{{}^cG,F}:=\mathcal R^c_{W_F,{}^cG}\times_{\mathcal R^c_{W_F,\mathbb G_m\times\Gamma_{\wt F/F}}}\{\chi\},\qquad
\Loc^{\wedge}_{{}^cG,F}=\Loc^{\wedge,\square}_{{}^cG,F}/\hatG^{\wedge}_\ell,
\end{equation*}
其中 $\hatG^{\wedge}_\ell$ 是 $\hatG$ 的 $\ell$-进完备化。由于 $\Gamma_F$ 是 $W_F$ 的投射有限完备化，引理 2.45 的一个轻微变体表明：$\Loc^{\wedge,\square}_{{}^cG,F}$ 在对应于 $\bar\rho:\Gamma_F\to{}^cG(\kappa)$ 的闭点处的完备化，就是由满足 $d\circ\rho=\chi$ 的 $\bar\rho$ 的加框架形变构成的空间 $\mathrm{Def}^{\square,\chi}_{\bar\rho}$。

\ysrc{71}
回忆 $\mathcal R^c_{W_F,{}^cG}$ 有到 $\Spec\Z_\ell$ 的推广 $\mathcal R^{sc}_{W_F,{}^cG}$，它分类 $W_F$ 的强连续表示（定义 2.56）。因此也可以把 (3.2) 推广到 $\Z_\ell$：
\begin{equation*}\tag{3.3}
\Loc^{\square}_{{}^cG,F}:=\mathcal R^{sc}_{W_F,{}^cG}\times_{\mathcal R^{sc}_{W_F,\mathbb G_m\times\Gamma_{\wt F/F}}}\{\chi\},\qquad
\Loc_{{}^cG,F}=\Loc^{\square}_{{}^cG,F}/\hatG_{\Z_\ell}.
\end{equation*}

\begin{rem}{3.2}
$\Spf\Z_p$ 上 $\Loc^{\wedge}_{{}^cG,F}$ 的类似物大概应当是 Emerton–Gee 叠 [23]（其定义复杂得多）。而 $\Spec\Z_p$ 上 $\Loc_{{}^cG,F}$ 的类似物则更微妙。
\end{rem}

\begin{rem}{3.3}
注意对 $\Loc^{\wedge}_{{}^cG,F}$ 而言，(2.36) 的分解就是按模 $\ell$ 惯性型所作的分解。事实上由 [7, 4.5]，(2.36) 中的 $\Theta$ 恰对应于 $I_F$ 的模 $\ell$ 完全可约表示（即模 $\ell$ 惯性型）。
\end{rem}

\begin{rem}{3.4}
由注 3.1，
\[
\Loc^{\square}_{{}^cG,F}\cong \mathcal R^{sc}_{W_F,{}^LG}\times_{\mathcal R^{sc}_{W_F,\Gamma_{\wt F/F}}}\{\mathrm{pr}\}
\]
在 $\Z_\ell[q^{\pm1/2}]$ 上成立。若 $G=\mathrm{GL}_m$，则 $\Loc^{\square}_{{}^cG,F}\cong\mathcal R^{sc}_{W_F,\mathrm{GL}_m}$。
\end{rem}

其次，有 $F$ 的 Weil–Deligne 表示的叠
\[
\Loc^{\mathrm{WD}}_{{}^cG,F}:=\Loc^{\mathrm{WD},\square}_{{}^cG,F}/\hatG,
\]
它是 $\mathbb Q$ 上的代数叠（见例如 [12, 2.1]）。这里 $\Loc^{\mathrm{WD},\square}_{{}^cG,F}$ 是按如下方式定义在 $\mathbf{CAlg}^{\heartsuit}_{\mathbb Q}$ 上的预层。记 $\hat{\mathcal N}_{\mathbb Q}\subset\mathrm{Lie}\,\hatG_{\mathbb Q}$ 为 $\hatG_{\mathbb Q}$ 的幂零锥。对 $\mathbb Q$-代数 $A$，给 ${}^cG(A)$ 赋离散拓扑，并令\tn{下述条件中的底本写作 $d\circ\rho=\chi$；此处表示记作 $r$，按上下文应为 $d\circ r=\chi$。}
\[
\Loc^{\mathrm{WD},\square}_{{}^cG,F}(A)=\Big\{(r,X)\ \Big|\ r:W_F\to{}^cG(A)\ \text{连续},\
X\in\hat{\mathcal N}_{\mathbb Q}(A)\ \Big|\ d\circ\rho=\chi,\ \mathrm{Ad}_{r(\gamma)}X=q^{\|\gamma\|}X\Big\}.
\]
注意 $\Loc^{\mathrm{WD},\square}_{{}^cG,F}$ 上有自然的 $\mathbb G_m$-作用，即对幂零元 $X$ 作伸缩。

可以看出
\[
\Loc^{\mathrm{WD},\square}_{{}^cG,F}=\varinjlim_L \Loc^{\mathrm{WD},\square}_{{}^cG,L/F},
\]
其中 $L$ 取遍 $F^{ur}\wt F$ 的所有在 $F$ 上伽罗瓦的有限扩张，而 $\Loc^{\mathrm{WD},\square}_{{}^cG,L/F}$ 是 $\Loc^{\mathrm{WD},\square}_{{}^cG,F}$ 的（开且闭的）子函子，由那些使 $r$ 穿过 $W_F/W_L\to{}^cG(A)$ 的 $(r,X)$ 组成。

由于 $W_F/W_L$ 是有限生成群，即 $\Z$ 被 $\Gamma_{L/F^{ur}}$ 的扩张，函子 $\Loc^{\mathrm{WD},\square}_{{}^cG,L/F}$ 可由 $\mathbb Q$ 上有限型的仿射概形表示。因此 $\Loc^{\mathrm{WD},\square}_{{}^cG,F}$ 与 $\Loc^{\mathrm{WD}}_{{}^cG,F}$ 是（ind-）可表示的。

\begin{rem}{3.5}
这里我们只把 $\Loc^{\mathrm{WD},\square}_{{}^cG,F}$ 定义为经典（ind-）概形，因为后续只需要这个。当然可以自然地把它定义为导出概形。但结果表明导出结构是平凡的。事实上，后面讨论 $\Loc^{\mathrm{WD},\square}_{{}^cG,F}$ 的整模型时我们会作这类讨论。
\end{rem}

\ysrc{72}
最后，在作出一个选择之后，我们可以把上述两个模空间粘成 $\Z[1/p]$ 上的代数叠。回忆以下基本事实（[42]）。

\begin{itemize}[leftmargin=1.6em,itemsep=0.2em]
\item 存在拓扑分裂 $\Gamma^t_F\to\Gamma_F$，使得 $\Gamma_F\cong P_F\rtimes\Gamma^t_F$。
\item 设 $\Gamma_q=\langle\tau,\sigma\rangle$ 如 (2.23)。则存在嵌入
\begin{equation*}\tag{3.4}
\iota:\Gamma_q\to\Gamma^t_F,
\end{equation*}
使得 $\iota(\tau)$ 是驯顺惯性的生成元，而 $\iota(\sigma)$ 是 Frobenius 的一个提升。于是 $\iota$ 诱导投影 $\Gamma_q\to\Z$ 的投射有限完备化与 $\Gamma^t_F\to\hat\Z$ 之间的同构。
\end{itemize}

对 $\iota$ 的一个选取，记 $\Gamma_{F,\iota}$ 为 $\Gamma_F$ 沿 $\iota$ 的拉回（不考虑这些群上的拓扑）。于是有包含 $\Gamma_{F,\iota}\to W_F\to\Gamma_F$。滥用记号，我们仍用 $\iota$ 同时记包含 $\Gamma_{F,\iota}\subset W_F$ 与 $\Gamma_{F,\iota}\subset\Gamma_F$。有短正合列
\[
1\to P_F\to\Gamma_{F,\iota}\to\Gamma_q\to1.
\]
(3.1) 中的同态 $\|\cdot\|$ 限制到 $\Gamma_{F,\iota}$。类似地，若 $L$ 在 $F^t$ 上有限且在 $F$ 上伽罗瓦，记 $\Gamma_{L/F,\iota}$ 为 $\Gamma_{L/F}$（即 $L/F$ 的伽罗瓦群）沿 $\iota$ 的拉回。有短正合列
\[
1\to Q_L:=\Gamma_{L/F^t}\to\Gamma_{L/F,\iota}\to\Gamma_q\to1,
\]
其中 $Q_L$ 是有限 $p$-群。

\begin{rem}{3.6}
(1) 注意对不同的选取 $\iota_1,\iota_2$，一般不存在限制在 $P_F$ 上为恒等映射的 $\Gamma_{F,\iota_1}$ 与 $\Gamma_{F,\iota_2}$ 之间的同构。

(2) (3.4) 中 $\iota$ 的所有可能选取构成 $\Aut^0$ 的一个主齐性空间（torsor），其中 $\Aut^0$ 是 $\Gamma^t_F$ 的连续自同构群中那些限制为 $I^t_F$ 的自同构、并在 $\Gamma^t_F/I^t_F$ 上诱导恒等映射的元素构成的群。$\Aut^0$ 本身是 $\hat\Z^{p,\times}:=\prod_{\ell\neq p}\Z_\ell^\times$ 被 $\hat\Z^p(1)$ 的扩张。
\end{rem}

现在选取 (3.4) 中的一个 $\iota$。若 $L/F^t\wt F$ 有限且 $L/F$ 伽罗瓦，则同态 $\chi_\iota:\Gamma_{F,\iota}\to\Z[1/p]^\times\times\Gamma_{\wt F/F}$ 穿过 $\Gamma_{L/F,\iota}\to\Z[1/p]^\times\times\Gamma_{\wt F/F}$，我们仍用同一记号记之，它可以看作 $\mathcal R_{\Gamma_{L/F,\iota},\mathbb G_m\times\Gamma_{\wt F/F}}$ 的一个 $\Z[1/p]$-点。定义概形
\[
\Loc^{\square}_{{}^cG,L/F,\iota}:=\mathcal R_{\Gamma_{L/F,\iota},{}^cG}\times_{\mathcal R_{\Gamma_{L/F,\iota},\mathbb G_m\times\Gamma_{\wt F/F}}}\{\chi_\iota\}.
\]
具体地说，对经典 $\Z[1/p]$-代数 $A$，
\[
\Loc^{\square}_{{}^cG,L/F,\iota}(A):=\Big\{\rho:\Gamma_{L/F,\iota}\to{}^cG(A)\ \Big|\ d\circ\rho=\chi_\iota:\Gamma_{L/F,\iota}\to\mathbb G_m\times\Gamma_{\wt F/F}\Big\}.
\]
现在把加框架 $\iota$-局部 Langlands 参数的概形定义为
\[
\Loc^{\square}_{{}^cG,F,\iota}:=\varinjlim_L \Loc^{\square}_{{}^cG,L/F,\iota}.
\]
同样由引理 2.45（的轻微变体），它在 $\bar\rho$ 处的形式完备化就是加框架形变空间 $\mathrm{Def}^{\square,\chi}_{\bar\rho}$。

\begin{prop}{3.7}
导出 ind-概形 $\Loc^{\square}_{{}^cG,F,\iota}$ 是有限型经典仿射概形的不交并，且在 $\Z[1/p]$ 上平坦。它是等维的，维数 $=\dim\hatG$，并且是局部完全交，其对偶复形平凡。
\end{prop}

\ysrc{73}
\begin{proof}
对 $\Gamma=\Gamma_{L/F,\iota}\simeq Q_L\rtimes\Gamma_q$ 以及 $M={}^cG$、$M=\mathbb G_m\times\Gamma_{\wt F/F}$ 应用命题 2.35。有投影 $\mathcal R_{\Gamma_{L/F,\iota},{}^cG}\to\mathcal R_{\Gamma_{L/F,\iota},\mathbb G_m\times\Gamma_{\wt F/F}}$。取 $\chi_\iota$ 上的纤维可知 $\Loc^{\square}_{{}^cG,L/F,\iota}$ 是有限型的经典仿射概形，在 $\Z[1/p]$ 上平坦，等维且维数 $=\dim\hatG$，并且是局部完全交。此外显然若 $L'/L$ 有限且 $L'/F$ 伽罗瓦，则 $\Loc^{\square}_{{}^cG,L/F,\iota}\subset\Loc^{\square}_{{}^cG,L'/F,\iota}$ 是开且闭的嵌入。命题得证。
\end{proof}

现在可以把 $\iota$-局部 Langlands 参数的叠定义为
\[
\Loc_{{}^cG,F,\iota}=\Loc^{\square}_{{}^cG,F,\iota}/\hatG.
\]
它是开且闭子叠 $\Loc_{{}^cG,L/F,\iota}=\Loc^{\square}_{{}^cG,L/F,\iota}/\hatG$ 的并，每个都在 $\Z[1/p]$ 上有限呈现。

\begin{rem}{3.8}
有两种方式把 $\Loc_{{}^cG,F,\iota}$（以及 $\Loc^{\mathrm{WD}}_{{}^cG,F}$）看作代数叠。第一种是把它看作局部有限型的叠，第二种是把它看作 ind-有限型叠。我们采用第二种观点。于是它的正则函数环（见下面 (3.5)）被看作 pro-代数。此外，后面我们会考虑 $\Loc_{{}^cG,F,\iota}$ 上凝聚层的范畴 $\Coh(\Loc_{{}^cG,F,\iota})$。按我们的定义，这些是只支撑在 $\Loc_{{}^cG,F,\iota}$ 的有限多个连通分支上的拟凝聚层复形，并且在这些分支上是凝聚复形。特别地，$\Loc_{{}^cG,F,\iota}$ 本身的结构层不被看作凝聚层，它属于 $\Coh(\Loc_{{}^cG,F,\iota})$ 的 ind-完备化 $\Ind\Coh(\Loc_{{}^cG,F,\iota})$。
\end{rem}

我们已经讨论了局部 Langlands 参数模的三个版本：一个在 $\Z_\ell$ 上，一个在 $\mathbb Q$ 上，一个在 $\Z[1/p]$ 上。下一项任务是把它们联系起来，并分析 $\Loc_{{}^cG,L/F,\iota}$ 对 $\iota$ 选取的依赖。

\begin{lem}{3.9}
映射 $\iota:\Gamma_{F,\iota}\to W_F$ 诱导自然同构\footnote{我们原先只考虑 $\Spf\Z_\ell$ 上的这种同构。感谢 P. Scholze 指出它在 $\Spec\Z_\ell$ 上也成立。}
\[
\phi_{\iota,\ell}:\Loc^{\square}_{{}^cG,F}\xrightarrow{\ \sim\ }\Loc^{\square}_{{}^cG,F,\iota}\otimes\Z_\ell.
\]
\end{lem}

\begin{proof}
先在经典模问题的层面证明。$\phi_{\iota,\ell}$ 把强连续表示 $W_F\to{}^cG(A)$ 送到它在 $\Gamma_{F,\iota}$ 上的限制。要证它是同构，只需证明每个 $\rho:\Gamma_{L/F,\iota}\to{}^cG(A)$ 都能延拓为 $W_F\to{}^cG(A)$ 的强连续表示。

如前，选取拓扑分裂 $\Gamma^t_F\to\Gamma_F$，写 $\Gamma_{F,\iota}\simeq P_F\rtimes\Gamma_q$。则存在某个 $N$（可能依赖于拓扑分裂的选取），使得 $\rho(\tau^N)\in\hat{\mathcal U}(A)$，其中 $\hat{\mathcal U}\subset\hatG$ 是 $\hatG$ 的幺幂簇。事实上，回忆子群 $\langle\tau\rangle\subset\Gamma_q$ 的限制诱导 $\Loc^{\square}_{{}^cG,F,\iota}\to{}^cG[q]$（见命题 2.33 的证明）。故只需证明存在 $N$，使得第 $N$ 次幂映射 ${}^cG\to{}^cG$，$g\mapsto g^N$ 把 ${}^cG[q]$ 送入 $\hat{\mathcal U}$。取忠实表示 ${}^cG\to\mathrm{GL}_m$，只需对 $\mathrm{GL}_m$ 证明类似陈述。这归结为：对 $X\in\mathrm{GL}_m$，若 $\mathrm{Char}(X^q)=\mathrm{Char}(X)$，则存在 $X$ 的某个幂 $X^N$ 使 $\mathrm{Char}(X^N)=(t-1)^m$。这是标准的。

\ysrc{74}
现在要证 $\rho$ 可以延拓，只需证明对每个 $X\in\hat{\mathcal U}(A)$，当 $A$ 赋以 ind-$\ell$-进拓扑时，映射 $\Z_\ell\to\hat{\mathcal U}(A)$，$a\mapsto X^a$ 连续。同样归约到 $\mathrm{GL}_m$ 情形。若 $\mathrm{Char}(X)=(t-1)^m$，则对每个 $v\in A^m$，$\{X^iv\}_{i\geq0}$ 含于一个有限 $\Z_\ell$-模。然后用引理 2.57 即可。

下面证明 $\phi_{\iota,\ell}$ 在导出层面也是同构。我们用命题 2.59 以及命题 2.42 中的论证，归约为证明
\[
C^*_{cts}(W_F,\mathrm{Ad}^0_\rho\otimes V)\to C^*(\Gamma_{F,\iota},\mathrm{Ad}^0_\rho\otimes V)
\]
对每个经典 $A$、每个通常 $A$-模 $V$、每个强连续同态 $\rho:W_F\to{}^cG(A)$ 都是同构。这里 $\mathrm{Ad}^0$ 是 ${}^cG$ 在 $\hatG$ 的李代数上的伴随表示。于是又归约为证明 $C^*_{cts}(I^t_F,(\mathrm{Ad}^0_\rho)^{P_F})\to C^*(\Z[1/p],(\mathrm{Ad}^0_\rho)^{P_F})$ 是同构。由引理 2.57，进一步归约为：若 $V$ 是 $I^t_F$ 在有限 $\Z_\ell$-模上的连续表示，则 $C^*_{cts}(I^t_F,V)\to C^*(\Z[1/p],V)$ 是同构。最后这个断言不难。
\end{proof}

另一方面，有下述结论。

\begin{lem}{3.10}
映射 $\Gamma_{F,\iota}\to W_F$ 诱导自然同构
\[
\phi_{\iota,\mathbb Q}:\Loc^{\mathrm{WD},\square}_{{}^cG,F}\xrightarrow{\ \sim\ }\Loc^{\square}_{{}^cG,F,\iota}\otimes\mathbb Q.
\]
\end{lem}

\begin{proof}
态射 $\phi_{\iota,\mathbb Q}$ 把 $(r,X)\in\Loc^{\mathrm{WD},\square}_{{}^cG,F}(A)$ 送到
\[
\rho:\Gamma_{F,\iota}\to{}^cG(A),\qquad \rho(\gamma)=r(\iota\gamma)\exp(|\gamma|_\iota X),
\]
其中 $|\gamma|_\iota\in\Z[1/p]$ 使得 $\gamma\in\Gamma_{F,\iota}$ 在 $\Gamma_q$ 中的像可写成 $\sigma^{\|\gamma\|}\tau^{|\gamma|_\iota}$，而
\[
\exp:\hat{\mathcal N}_{\mathbb Q}\xrightarrow{\ \sim\ }\hat{\mathcal U}_{\mathbb Q}
\]
是通常的指数映射，它在 $\hatG$ 的幂零簇与幺幂簇之间（在 $\mathbb Q$ 上）诱导同构。记 $\log:\hat{\mathcal U}_{\mathbb Q}\to\hat{\mathcal N}_{\mathbb Q}$ 为它的逆。

下面定义反方向的态射。设 $\rho:\Gamma_{F,\iota}\to{}^cG(A)$ 是 $\Loc^{\square}_{{}^cG,F,\iota}$ 的一个 $A$-点，并设它穿过某个 $\Gamma_{L/F,\iota}$。注意存在 $m$，使得 $\tau^m\in\Gamma_q$ 在 $\Gamma_{F,\iota}$ 中的像不依赖于分裂 $\Gamma_q\to\Gamma_{L/F,\iota}$ 的选取。此外，把 $m$ 换成它的倍数，可设 $\rho(\tau)^m\in\hat{\mathcal U}_{\mathbb Q}(A)$。于是取 $X=\frac1m\log(\rho(\tau)^m)$。显然 $X$ 不依赖 $m$ 的选取。于是得到良定义的同态
\[
r:\Gamma_{F,\iota}\to{}^cG(A),\qquad r(\gamma)=\rho(\gamma)\exp(-|\gamma|_\iota X).
\]
由于 $r(\tau^m)=1$，可以把 $r$ 看作连续映射 $W_{L/F}\to{}^cG(A)$，其中 $A$ 赋以离散拓扑。于是 $\rho\mapsto(r,X)$ 给出 $\phi_{\iota,\mathbb Q}$ 的逆。
\end{proof}

继续之前，注意作为副产品我们得到以下结论。

\begin{cor}{3.11}
概形 $\Loc^{\square}_{{}^cG,F,\iota}$ 是约化的。
\end{cor}

注意 $\Loc^{\square}_{{}^cG,F,\iota}$ 在某个素数 $\ell$ 上的纤维可能非约化。

\begin{proof}
由于 $\Loc^{\square}_{{}^cG,F,\iota}$ 是在 $\Z[1/p]$ 上平坦的局部完全交（命题 3.7），结论由 $\Loc^{\square}_{{}^cG,F,\iota}\otimes\mathbb Q\cong\Loc^{\mathrm{WD},\square}_{{}^cG,F}$ 的一般光滑性（generic smoothness，已在 [12] 中证明）与 Serre 判别法 S1 得到。
\end{proof}

\ysrc{75}
现在可以比较不同选取 $\iota$ 下的 $\Loc^{\square}_{{}^cG,F,\iota}$。设 $\iota_1,\iota_2:\Gamma_q\to\Gamma^t_F$ 是两个嵌入。由注 3.6 回忆，存在 $\vartheta\in\Aut^0$ 使 $\iota_2=\vartheta\iota_1:\Gamma_q\to\Gamma^t_F$，且有投影 $\Aut^0\to\Z_\ell^\times$。记 $\bar\vartheta\in\Z_\ell^\times$ 为 $\vartheta$ 的像。由于 $\mathbb G_m$ 通过伸缩幂零元作用在 $\Loc^{\mathrm{WD},\square}_{{}^cG,F}$ 上，$\bar\vartheta$（看作 $\mathbb G_m(\mathbb Q_\ell)$ 中的元素）作用在 $\Loc^{\mathrm{WD},\square}_{{}^cG,F}\otimes\mathbb Q_\ell$ 上。

\begin{prop}{3.12}
存在唯一的同构 $\vartheta=\vartheta_{\iota_1,\iota_2}:\Loc^{\square}_{{}^cG,F,\iota_1}\otimes\Z_\ell\cong\Loc^{\square}_{{}^cG,F,\iota_2}\otimes\Z_\ell$（$\Z_\ell$ 上的概形），使下图交换
\begin{center}
\begin{tikzcd}[column sep=large,row sep=large]
\Loc^{\square}_{{}^cG,F} \arrow[r,"\phi_{\iota_1,\ell}"] \arrow[d,equal]
& \Loc^{\square}_{{}^cG,F,\iota_1}\otimes\Z_\ell \arrow[d,"\vartheta"]
& \Loc^{\mathrm{WD},\square}_{{}^cG,F}\otimes\Q_\ell \arrow[l,"\phi_{\iota_1,\Q_\ell}"'] \arrow[d,"\bar\vartheta"] \\
\Loc^{\square}_{{}^cG,F} \arrow[r,"\phi_{\iota_2,\ell}"']
& \Loc^{\square}_{{}^cG,F,\iota_2}\otimes\Z_\ell
& \Loc^{\mathrm{WD},\square}_{{}^cG,F}\otimes\Q_\ell \arrow[l,"\phi_{\iota_2,\Q_\ell}"]
\end{tikzcd}
\end{center}
\end{prop}

\begin{proof}
由于 $\phi_{\iota_i,\ell}$ 是同构，故存在唯一的与诸 $\phi_{\iota_i,\ell}$ 相容的 $\vartheta$。追踪构造可见 $\vartheta\circ\phi_{\iota_1,\mathbb Q_\ell}=\phi_{\iota_2,\mathbb Q_\ell}\circ\bar\vartheta$。
\end{proof}

\begin{cor}{3.13}
$\Loc_{{}^cG,F,\iota}$ 上的正则函数环
\begin{equation*}\tag{3.5}
Z_{{}^cG,F}:=H^0\Gamma(\Loc_{{}^cG,F,\iota},\mathcal O)
\end{equation*}
不依赖 $\iota$ 的选取，至多相差一个典范同构（故可略去下标 $\iota$）。
\end{cor}

回忆按我们的约定，$\Gamma(\Loc_{{}^cG,F,\iota},-)$ 表示导出函子，而 $H^0\Gamma$ 表示它的零次上同调。

\begin{proof}
事实上，$\mathbb G_m$ 在 $\Loc^{\mathrm{WD}}_{{}^cG,F}$ 上的作用（伸缩幂零元）在其正则函数环上诱导平凡作用。因此命题 3.12 中的 $\bar\vartheta$ 在取 $\hatG$-不变量之后诱导恒等映射。
\end{proof}

这个代数通常称为 $G^*$（$G$ 的拟分裂内型）的稳定中心（stable center），至少在基变换到 $\mathbb C$ 时如此（见 [37]）。它带有按 $\Loc_{{}^cG,F,\iota}$ 的连通分支指标化的幂等分解。对连通分支的有限并 $D$，记 $Z_{{}^cG,F,D}$ 为相应的正则函数环，它是有限生成 $\Lambda$-代数。若 $D=\Loc_{{}^cG,L/F,\iota}$，记 $Z_{{}^cG,F,D}$ 为 $Z_{{}^cG,L/F}$。

由于当 $\Lambda$ 不是特征零的域时，在 $\Lambda$ 上取 $\hatG$-不变量对 $\hatG$-表示不必正合，先验地高阶上同调 $H^i\Gamma(\Loc_{{}^cG,F,\iota},\mathcal O)$ 对 $i>0$ 可能不消失。但猜想 4.24 提示并非如此。事实上我们提出如下猜想\footnote{事实上，当 $\ell$ 不太小时，这一猜想已在 [27] 中得证。}。

\begin{conj}{3.14}
对每个 $i\geq1$，$H^i\Gamma(\Loc_{{}^cG,F,\iota},\mathcal O)=0$。
\end{conj}

\begin{rem}{3.15}
设 $\kappa$ 是 $\Z[1/p]$ 上的代数闭域。由 [49, 11.7] 与 [7, 4.5] 以及注 2.26，$Z_{{}^cG,F}$ 的 $\kappa$-点与满足以下条件的同态 $\rho:\Gamma_{F,\iota}\to{}^cG(\kappa)$ 的 $\hatG(\kappa)$-共轭类之间存在双射：
\begin{itemize}[leftmargin=1.6em,itemsep=0.2em]
\item $d\circ\rho=\chi$；
\item $\rho$ 穿过 $\Gamma_{L/F,\iota}\to{}^cG(\kappa)$，其中 $L/F^t\wt F$ 是某个有限扩张；
\item $\rho$ 完全可约（按 [7, 3.5] 的意义）。
\end{itemize}
\end{rem}

\ysrc{76}
给定猜想 3.14，还可以进一步猜测 (2.21) 在当前设定下的一个轻微变体是同构（取 $\pi_0$ 之后）。

本小节最后讨论这些叠在张量诱导下的行为。设 $F'/F$ 是有限可分扩张。设 $G'$ 是 $F'$ 上的连通约化群，$G=\mathrm{Res}_{F'/F}G'$。如 [9, 5.1,4.1] 所述，$G$ 的对偶群 $\hatG$（带 $\Gamma_F$-作用）典范同构于张量诱导 $\mathrm{Ind}^{\Gamma_F}_{\Gamma_{F'}}\hatG'$，后者按定义是从 $\Gamma_F$ 到 $\hatG'$ 的所有 $\Gamma_{F'}$-等变映射构成的空间。有 $\Gamma_{F'}$-等变映射（[9, 4.1]）
\[
\hatG'\xrightarrow{i}\hatG\xrightarrow{\mathrm{ev}_e}\hatG',
\]
其复合为恒等，其中第一个映射把 $g$ 送到唯一的映射 $f:\Gamma_F\to\hatG'$：它支撑在 $\Gamma_{F'}$ 上且 $f(1)=g$；第二个映射把 $f:\Gamma_F\to\hatG'$ 送到 $f(e)$。于是有典范同态 ${}^c(G')\to{}^cG$，它与 $i$ 以及 $\mathbb G_m\times\Gamma_{\wt F'/F'}\to\mathbb G_m\times\Gamma_{\wt F/F}$ 相容，如 [9, 5.1 (5)] 所述。$\iota:\Gamma_q\to\Gamma^t_F$ 的选取给出 $\iota':\Gamma_{q'}\to\Gamma^t_{F'}$。注意 $\mathrm{Ind}^{\Gamma_{F,\iota}}_{\Gamma_{F',\iota'}}\hatG'=\mathrm{Ind}^{\Gamma_F}_{\Gamma_{F'}}\hatG'$。

\begin{lem}{3.16}
有典范同构
\[
\Loc_{{}^cG,F,\iota}\cong\Loc_{{}^cG',F',\iota'},\qquad \rho\mapsto \mathrm{ev}_e\circ(\rho|_{\Gamma_{F',\iota'}}).
\]
\end{lem}

\begin{proof}
这是 Shapiro 引理的几何版本。我们把 [78, 4.1.2] 的论证推广，以显式构造逆映射。为简单起见，写 $\Gamma'=\Gamma_{F',\iota'}$，$\Gamma=\Gamma_{F,\iota}$。设 $s:\Gamma'\backslash\Gamma\to\Gamma$ 是投影 $\Gamma\to\Gamma'\backslash\Gamma$，$\gamma\mapsto\bar\gamma$ 的一个截面（它把单位陪集送到 $1\in\Gamma$）。于是有映射
\[
\Xi_s:\Gamma\to\Gamma',\qquad \Xi_s(\gamma):=\gamma s_{\bar\gamma}^{-1}.
\]
注意对 $\gamma'\in\Gamma'$ 有 $\Xi_s(\gamma'\gamma)=\gamma'\Xi_s(\gamma)$。此外令
\[
\Delta_s:\hatG'\to\hatG,\qquad \Delta_s(g):\Gamma\to\hatG',\qquad \Delta_s(g)(\delta)=\chi(\Xi_s(\delta))(g).
\]
现在如下构造态射 $I_s:\Loc^{\square}_{{}^cG',F',\iota'}\to\Loc^{\square}_{{}^cG,F,\iota}$。设 $\rho'=(\phi',\chi):\Gamma'\to{}^c(G')(A)=\hatG'(A)\rtimes(A^\times\times\Gamma_{\wt F'/\wt F'})$。定义 $I_s(\rho')=(\phi,\chi):\Gamma\to{}^cG(A)=\hatG(A)\rtimes(A^\times\times\Gamma_{\wt F/\wt F})$，其中
\[
\phi(\gamma):\Gamma\to\hatG'(A),\qquad \phi(\gamma)(\delta)=\phi'(\Xi_s(\delta))^{-1}\phi'(\Xi_s(\delta\gamma)).
\]
可以验证
\begin{itemize}[leftmargin=1.6em,itemsep=0.2em]
\item 对 $\gamma'\in\Gamma'$ 有 $\phi(\gamma'\gamma)=\chi(\gamma')(\phi(\gamma))$，故 $\phi(\gamma)\in\hatG(A)$；
\item $I_s(\rho')$ 是同态 $\Gamma\to{}^cG(A)$，且 $\mathrm{ev}_e\circ(I_s(\rho')|_{\Gamma'})=\rho'$；
\item 对任意 $g\in\hatG'(A)$ 有 $I_s(g^{-1}\rho'g)=\Delta_s(g)^{-1}I_s(\rho')\Delta_s(g)$。
\end{itemize}
于是我们得到与引理中映射互逆的态射 $\Loc_{{}^cG',F',\iota'}\to\Loc_{{}^cG,F,\iota}$。
\end{proof}

\subsection{环面的对偶性与 $\Coh(\Loc_{{}^cG,F,\iota})$ 的对称性}

先考察 $G=T$ 是 $F$ 上环面的情形下的叠 $\Loc_{{}^cG,F,\iota}$。由命题 3.12 的证明不难看出 $\Loc_{{}^cT,F,\iota}$ 不依赖 $\iota$ 的选取。但事实上可以如下显式描述 $\Loc_{{}^cT,F,\iota}$。设 $\wt F/F$ 是 $T$ 的分裂域。由局部类域论，有短正合列
\[
1\to\wt F^\times\to W_{\wt F/F}\to\Gamma_{\wt F/F}\to1,
\]

\ysrc{77}
其中 $W_{\wt F/F}$ 是扩张 $\wt F/F$ 的 Weil 群。设 $U^{(n)}$ 为 $\wt F$ 的第 $n$ 个单位群（故 $U^{(0)}=\mathcal O_{\wt F}^\times$，且 $n\geq1$ 时 $U^{(n)}=1+\mathfrak m_{\wt F}^n$），并写 $W^{(n)}=W_{\wt F/F}/U^{(n)}$。则有自然同构
\[
\Loc_{{}^cT,F,\iota}\cong\varinjlim_n \Loc^{(n)}_{{}^cT,F},
\]
其中
\[
\Loc^{(n)}_{{}^cT,F}:={}^{cl}\Big(\mathcal R^{\square}_{W^{(n)},{}^cT}\times_{\mathcal R^{\square}_{W^{(n)},\mathbb G_m\times\Gamma_{\wt F/F}}}\{\chi\}\Big)/\widehat{T}.
\]
于是今后略去记号中的下标 $\iota$。

\begin{exam}{3.17}
设 $\wt F/F$ 驯顺分歧，故 $W^{(1)}$ 是 $\Gamma^t_F$ 的商。则 $\Loc^{(1)}_{{}^cT,F}$ 就是下文将引入的 $T$ 的驯顺 Langlands 参数叠 $\Loc^{\mathrm{tame}}_{{}^cT,F}$。注意 $\Loc^{(1)}_{{}^cT,F}$ 在 $\Z[1/p]$ 上连通，但在 $\mathbb Q$ 上不连通。

若 $\wt F/F$ 非分歧，则 $\Loc^{(0)}_{{}^cT,F}$ 也可等同于下文将引入的 $T$ 的非分歧参数叠 $\Loc^{\mathrm{ur}}_{{}^cT,F}$。此时记 $\bar\sigma$ 为 $\Gamma_{\wt F/F}$ 中的 Frobenius 元。则包含 $\Loc^{(0)}_{{}^cT,F}\subset\Loc^{(1)}_{{}^cT,F}$ 等同于
\begin{equation*}\tag{3.6}
\widehat{T}^{\bar\sigma}/\widehat{T}=\{1\}\times\widehat{T}^{\bar\sigma}/\widehat{T}
\subset\Big(\big({}^{cl}\mathcal R_{\kappa_{\wt F}^\times,\widehat{T}}\big)^\sigma\times\widehat{T}^{\bar\sigma}\Big)/\widehat{T}.
\end{equation*}
这里 $\big({}^{cl}\mathcal R_{\kappa_{\wt F}^\times,\widehat{T}}\big)^\sigma$ 是从 $\kappa_{\wt F}^\times$ 到 $\widehat{T}$ 的 $\sigma$-等变同态构成的经典模空间，$1$ 表示平凡同态。（注意如例 2.30 所解释，$\mathcal R_{\kappa_{\wt F}^\times,\widehat{T}}$ 本身并不经典（当 $\ell\mid\#\kappa_{\wt F}^\times-1$ 时在 $\mathbb F_\ell$ 上并不经典），故需取其底层的经典概形。）
\end{exam}

注意 $\Loc_{{}^cT,F}$ 实际上是 $\Z[1/p]$ 上的 Picard 叠（关于 Picard 叠的一般综述见例如 [15, §A]）。设 $\mathbb{B}\mathbb G_m$ 为 $\Z[1/p]$ 上 $\mathbb G_m$ 的分类叠。令
\[
\Loc^{\vee}_{{}^cT,F}:=\Hom(\Loc_{{}^cT,F},\mathbb{B}\mathbb G_m)
\]
为 $\Loc_{{}^cT,F}$ 在 $\Z[1/p]$ 上的对偶 Picard 叠（[15, A.3.1] 意义下），它仍是 Picard 叠，分类 $\Loc_{{}^cT,F}$ 上的乘性线丛。另一方面，设 $\breve F$ 为 $F$ 的极大非分歧扩张 $F^{ur}/F$ 的完备化。则 Frobenius $\sigma$ 作用在 $\breve F$ 上。记 $\mathbf{Tor}_{T,\mathrm{iso}F}$ 为由 $(\mathcal E,\varphi)$ 构成的 Picard 群胚，其中 $\mathcal E$ 是 $\breve F$ 上的 $T$-torsor，$\varphi:\mathcal E\simeq\sigma^*\mathcal E$ 是 $T$-torsor 的同构。（可以把 $(\mathcal E,\varphi)$ 看作 [46, §3] 与 [47, §2] 意义下 $F$-线性的 $\sigma$-$\breve F$-空间 Tannaka 范畴中的 $T$-torsor。）把 $\mathbf{Tor}_{T,\mathrm{iso}F}$ 看作 $\Z[1/p]$ 上的常 Picard 叠。下面的猜想可以看作非阿基米德局部域上环面的局部 Langlands 对偶。

\begin{conj}{3.18}
存在 $\mathbf{Tor}_{T,\mathrm{iso}F}\times\Loc_{{}^cT,F}$ 上的自然 Poincaré 线丛，它诱导 Picard 叠的同构 $\mathbf{Tor}_{T,\mathrm{iso}F}\cong\Loc^{\vee}_{{}^cT,F}$。
\end{conj}

\begin{rem}{3.19}
注意 $\mathbf{Tor}_{T,\mathrm{iso}F}$ 的同构类不是别的，正是 $T$ 的 Kottwitz 集 $B(T)$（见 [46,47]），在 loc.\ cit.\ 中它等同于 $X_*(\widehat{T}^{\Gamma_F})$。另一方面，每个 $T$-torsor 的自同构群就是 $T(F)$，其特征标群可以等同于 $T$ 的 Langlands 参数集（[52]）。故该猜想是这些事实的代数几何精细化。
\end{rem}

我们把这个猜想稍作推广，允许 $F$ 上不必连通的乘法型群 $Z$。Picard 群胚 $\mathbf{Tor}_{Z,\mathrm{iso}F}$ 仍有意义（如 [47]），但现在可能有非平凡的导出结构（因为 $H^2(W_F,Z(\bar F))$
